\chapter{From an order-seven automorphism to a canonical one}\label{sec:group}

This chapter carries out the group-theoretic reduction: the statement
``$7 \nmid |\Aut\Gamma|$ for every $\srg(99,14,1,2)$ $\Gamma$'' is reduced to a
single concrete assertion about one explicit permutation $\sigma_0$ of
$\{0,\dots,98\}$, namely that no $\sigma_0$-invariant $\srg(99,14,1,2)$ exists
(Theorem~\ref{thm:reduction}).  The mathematical content follows
Cesarz--Woldar: an order-seven automorphism of an $\srg(99,14,1,2)$ fixes
exactly one vertex, hence has cycle type $[1,7^{14}]$, hence is conjugate in
the symmetric group to the canonical block rotation $\sigma_0$.

\section{Fixed points of prime-order permutations}

We begin with elementary counting lemmas about permutations $\pi$ with
$\pi^p = 1$ for a prime $p$.  Throughout, $\Fix\pi$ denotes the set of fixed
points and $\operatorname{supp}\pi$ its complement, the support.

\begin{lemma}[Support is a multiple of $p$]\label{lem:grp-support-multiple}
Let $\pi$ be a permutation of a finite type with $\pi^p = 1$ for a prime $p$.
Then $|\operatorname{supp}\pi|$ is a multiple of $p$.
\lean{ConwayO7.card_support_of_pow_prime_eq_one}
\leanfile{ConwayO7/CycleType.lean}\leanok
\end{lemma}

\begin{proof}\leanok
Either $\pi = 1$, with empty support, or $\ord\pi = p$; in the latter case the
cycle type of $\pi$ is a multiset of copies of $p$, and the support is the
disjoint union of the cycles.
\end{proof}

\begin{lemma}[Fixed points and support partition the domain]\label{lem:grp-fix-support}
For any permutation $\pi$ of a finite type $\alpha$,
$|\Fix\pi| + |\operatorname{supp}\pi| = |\alpha|$.
\lean{ConwayO7.card_filter_fixed_add_card_support}
\leanfile{ConwayO7/CycleType.lean}\leanok
\end{lemma}

\begin{lemma}[Fixed-point congruence]\label{lem:grp-fixed-congruence}
Let $\pi$ be a permutation with $\pi^p = 1$, $p$ prime, and let $s$ be a
finite $\pi$-invariant set.  Then
\[
  |\{a \in s \mid \pi a = a\}| \equiv |s| \pmod p .
\]
\lean{ConwayO7.card_filter_fixed_modEq}
\leanfile{ConwayO7/CycleType.lean}\leanok
\end{lemma}

\begin{proof}\leanok
\uses{lem:grp-support-multiple, lem:grp-fix-support}
Since $\pi$ is injective and $s$ is finite, $\pi(s) \subseteq s$ upgrades to
$\pi(s) = s$, so $\pi$ restricts to a permutation $\pi'$ of $s$ with
$(\pi')^p = 1$.  By Lemmas~\ref{lem:grp-support-multiple}
and~\ref{lem:grp-fix-support} applied to $\pi'$, the fixed points of $\pi$
inside $s$ number $|s| - pm$ for some $m$.  (We will apply this with $p = 7$
to the vertex set and to a neighbourhood, and with $p = 2$ to a matching
involution.)
\end{proof}

\begin{proposition}[Cycle type from a fixed-point count]\label{lem:grp-cycletype-of-fixed}
Let $\pi \ne 1$ be a permutation of a finite type of cardinality $f + pm$ with
$\pi^p = 1$, $p$ prime, and exactly $f$ fixed points.  Then the cycle type of
$\pi$ is $[p^m]$.
\lean{ConwayO7.cycleType_eq_replicate_of_fixed}
\leanfile{ConwayO7/Canonical.lean}\leanok
\end{proposition}

\begin{proof}\leanok
\uses{lem:grp-fix-support}
Since $\pi \ne 1$ and $\pi^p = 1$, the order of $\pi$ is exactly $p$, so its
cycle type is $[p^n]$ for some $n \ge 1$.  By
Lemma~\ref{lem:grp-fix-support} the support has $pm$ elements, and the cycle
lengths sum to the support size, so $pn = pm$ and $n = m$.
\end{proof}

\section{Local structure of an \texorpdfstring{$\srg(99,14,1,2)$}{srg(99,14,1,2)}}

Fix an $\srg(99,14,1,2)$ $\Gamma$ on a vertex set $V$, and write
$N(x)$ for the neighbourhood of $x$ and $N(x,y) = N(x) \cap N(y)$ for a
common neighbourhood.

\begin{lemma}[$\lambda = 1$: unique triangles]\label{lem:grp-unique-triangle}
If $x$ and $u$ are adjacent then $N(x,u)$ is a singleton: every edge of
$\Gamma$ lies in exactly one triangle.
\lean{ConwayO7.commonNeighbors_eq_singleton_of_adj}
\leanfile{ConwayO7/CycleType.lean}\leanok
\uses{def:srg}
\end{lemma}

\begin{lemma}[$\mu = 2$: common-neighbour pairs]\label{lem:grp-mu-pair}
If $x \ne w$ are nonadjacent then $N(x,w) = \{a,b\}$ for two distinct
vertices $a \ne b$.
\lean{ConwayO7.commonNeighbors_eq_pair_of_not_adj}
\leanfile{ConwayO7/CycleType.lean}\leanok
\uses{def:srg}
\end{lemma}

\begin{lemma}[Injectivity of the $\Gamma_2$-labelling]\label{lem:grp-label-injective}
Let $w, w' \ne x$ be nonadjacent to $x$.  If $N(x,w) = N(x,w')$ then
$w = w'$: a vertex at distance two from $x$ is determined by its pair of
common neighbours with $x$.
\lean{ConwayO7.eq_of_commonNeighbors_eq}
\leanfile{ConwayO7/CycleType.lean}\leanok
\uses{def:srg}
\end{lemma}

\begin{proof}\leanok
\uses{lem:grp-unique-triangle, lem:grp-mu-pair}
Write $N(x,w) = \{a,b\}$ with $a \ne b$.  The labels $a, b$ are nonadjacent:
an edge $ab$ would put the two distinct vertices $x$ and $w$ into the unique
triangle over $ab$ granted by $\lambda = 1$.  Hence by $\mu = 2$ the pair
$\{a,b\}$ has exactly two common neighbours; both $\{x,w\}$ and $\{x,w'\}$
are contained in $N(a,b)$ and have the right size, so
$\{x,w\} = N(a,b) = \{x,w'\}$ and $w = w'$.
\end{proof}

\begin{lemma}[Equivariance of common neighbourhoods]\label{lem:grp-equivariance}
Let $\pi$ be a permutation of $V$ with
$\Gamma(\pi u, \pi v) \Leftrightarrow \Gamma(u,v)$ for all $u,v$.  Then
$\pi\bigl(N(v,w)\bigr) \subseteq N(\pi v, \pi w)$.
\lean{ConwayO7.mapsTo_commonNeighbors}
\leanfile{ConwayO7/CycleType.lean}\leanok
\end{lemma}

\begin{proposition}[Rigidity]\label{lem:grp-rigidity}
An adjacency-preserving permutation of an $\srg(99,14,1,2)$ that fixes a
vertex $x$ and every neighbour of $x$ is the identity.
\lean{ConwayO7.eq_one_of_fixes_neighbors}
\leanfile{ConwayO7/CycleType.lean}\leanok
\uses{def:srg}
\end{proposition}

\begin{proof}\leanok
\uses{lem:grp-label-injective, lem:grp-equivariance}
Let $\pi$ fix $x$ and $N(x)$ pointwise, and let $v \notin \{x\} \cup N(x)$.
The common neighbourhood $N(x,v)$ consists of neighbours of $x$, hence is
pointwise fixed; by equivariance (applied to $\pi$ and to $\pi^{-1}$) it
follows that $N(x, \pi v) = N(x, v)$.  Since $\pi v$ is again nonadjacent to
$x$ and distinct from $x$, injectivity of the labelling
(Lemma~\ref{lem:grp-label-injective}) forces $\pi v = v$.
\end{proof}

\section{The fixed-point count and the cycle type}

\begin{theorem}[Exactly one fixed vertex]\label{lem:fixed}
Let $\Gamma$ be an $\srg(99,14,1,2)$ and let $\pi$ be an adjacency-preserving
permutation of its vertices with $\ord\pi = 7$.  Then $\pi$ fixes exactly one
vertex.
\lean{ConwayO7.card_filter_fixed_eq_one}
\leanfile{ConwayO7/CycleType.lean}\leanok
\uses{def:srg}
\end{theorem}

\begin{proof}[Proof sketch]\leanok
\uses{lem:grp-fixed-congruence, lem:grp-unique-triangle, lem:grp-mu-pair,
lem:grp-equivariance, lem:grp-rigidity}
\emph{Step 0.}  By the congruence (Lemma~\ref{lem:grp-fixed-congruence},
$p = 7$, $s = V$), $|\Fix\pi| \equiv 99 \equiv 1 \pmod 7$; in particular a
fixed vertex $x$ exists.

\emph{Step 1.}  Let $s = N(x) \cap \Fix\pi$ be the set of fixed neighbours of
$x$.  By $\lambda = 1$, the map sending $u \in N(x)$ to the apex of the
unique triangle over the edge $xu$ is a fixed-point-free involution of
$N(x)$; by equivariance and uniqueness of the apex it preserves $s$, so
(Lemma~\ref{lem:grp-fixed-congruence} with $p = 2$) $|s|$ is even.  Applying
the congruence with $p = 7$ to the invariant set $N(x)$ gives
$|s| \equiv 14 \equiv 0 \pmod 7$; together with $|s| \le 14$ this leaves
$|s| \in \{0, 14\}$.

\emph{Step 2.}  If $|s| = 14$ then $\pi$ fixes $x$ and all of $N(x)$, so
$\pi = 1$ by rigidity (Proposition~\ref{lem:grp-rigidity}), contradicting
$\ord\pi = 7$.  Hence $s = \emptyset$.

\emph{Step 3.}  Suppose $w \ne x$ were a second fixed vertex.  It cannot be
adjacent to $x$ (it would lie in $s$).  Otherwise $\pi$ permutes the two
common neighbours $\{a,b\} = N(x,w)$; neither is fixed (both are neighbours
of $x$), so $\pi$ swaps them.  Then $\pi^2$ fixes $a$, and
$\pi = (\pi^2)^4$ fixes $a \in N(x)$ --- contradiction.  So $\Fix\pi = \{x\}$.
\end{proof}

\begin{corollary}[Cycle type of an order-seven automorphism]\label{cor:grp-aut-cycle-type}
Let $\Gamma$ be an $\srg(99,14,1,2)$ and $\sigma \in \Aut\Gamma$ with
$\ord\sigma = 7$.  Then the underlying vertex permutation of $\sigma$ has
cycle type $[7^{14}]$.
\lean{ConwayO7.cycleType_autToPerm}
\leanfile{ConwayO7/CycleType.lean}\leanok
\uses{def:srg, def:grp-autToPerm, thm:grp-cycle-type}
\end{corollary}

\begin{proof}\leanok
\uses{lem:grp-autToPerm-injective}
An injective homomorphism preserves orders, so the underlying permutation
also has order $7$ and preserves adjacency; apply
Theorem~\ref{thm:grp-cycle-type}.
\end{proof}

\section{The canonical permutation \texorpdfstring{$\sigma_0$}{sigma0}}

\begin{definition}[The canonical permutation]\label{def:sigma0}
On the vertex set $\{0,\dots,98\}$, partition the nonzero vertices into
fourteen consecutive blocks $B_i = \{1+7i, \dots, 7+7i\}$, $0 \le i \le 13$.
Let $\sigma_0$ be the permutation fixing $0$ and rotating each block
cyclically:
\[
  \sigma_0(0) = 0, \qquad
  \sigma_0(1 + 7i + j) = 1 + 7i + \bigl((j+1) \bmod 7\bigr)
  \quad (0 \le j < 7).
\]
Formally, $\sigma_0$ is assembled from the explicit forward-rotation formula
and its backward-rotation companion (positions advance by $1$, respectively
by $6 \equiv -1$, modulo $7$).
\lean{ConwayO7.sigma0, ConwayO7.sigma0Fun, ConwayO7.sigma0InvFun,
ConwayO7.sigma0_val}
\leanfile{ConwayO7/Canonical.lean}\leanok
\uses{lem:grp-sigma0-inverse}
\end{definition}

\begin{proof}\leanok
On a nonzero vertex $1+7i+j$ the composite advances the position by
$1 + 6 \equiv 0 \pmod 7$ and leaves the block index $i$ unchanged; the vertex
$0$ is fixed by both maps.
\end{proof}

\begin{lemma}[Blocks are constant, positions advance]\label{lem:grp-sigma0-iterate}
For a nonzero vertex $v = 1 + 7i + j$ and every $k \ge 0$,
\[
  \sigma_0^{\,k}(v) = 1 + 7i + \bigl((j+k) \bmod 7\bigr):
\]
iterating $\sigma_0$ never leaves the block of $v$ and advances the position
within the block by $k$ modulo $7$.
\lean{ConwayO7.sigma0_iterate_val}
\leanfile{ConwayO7/Canonical.lean}\leanok
\uses{def:sigma0}
\end{lemma}

\begin{proof}\leanok
Induction on $k$: a nonzero vertex stays nonzero (its position formula keeps
it in the block $B_i$), so the defining formula of $\sigma_0$ applies at every
step and adds $1$ to the position, modulo $7$.
\end{proof}

\begin{lemma}[$\sigma_0$ has order dividing seven]\label{lem:grp-sigma0-order}
$\sigma_0^7 = 1$.
\lean{ConwayO7.sigma0_pow_seven, ConwayO7.sigma0_iterate_seven}
\leanfile{ConwayO7/Canonical.lean}\leanok
\uses{def:sigma0}
\end{lemma}

\begin{proof}\leanok
\uses{lem:grp-sigma0-iterate}
The vertex $0$ is fixed; on a nonzero vertex,
Lemma~\ref{lem:grp-sigma0-iterate} with $k = 7$ advances the position by
$7 \equiv 0 \pmod 7$.
\end{proof}

\begin{lemma}[Fixed points of $\sigma_0$]\label{lem:grp-sigma0-fixed-iff}
$\sigma_0(v) = v$ if and only if $v = 0$.
\lean{ConwayO7.sigma0_fixed_iff}
\leanfile{ConwayO7/Canonical.lean}\leanok
\uses{def:sigma0}
\end{lemma}

\begin{proof}\leanok
On a nonzero vertex the position advances, and $(j+1) \bmod 7 = j$ is
impossible for $0 \le j < 7$.
\end{proof}

\begin{lemma}[Fixed-point count of $\sigma_0$]\label{lem:grp-sigma0-fixed-card}
$\sigma_0$ has exactly one fixed point; in particular $\sigma_0 \ne 1$.
\lean{ConwayO7.sigma0_fixed_card, ConwayO7.sigma0_ne_one}
\leanfile{ConwayO7/Canonical.lean}\leanok
\uses{lem:grp-sigma0-fixed-iff}
\end{lemma}

\begin{lemma}[Cycle type of $\sigma_0$]\label{lem:grp-sigma0-cycle-type}
The cycle type of $\sigma_0$ is $[7^{14}]$ --- the same as that of any
order-seven automorphism of an $\srg(99,14,1,2)$
(Corollary~\ref{cor:grp-aut-cycle-type}).
\lean{ConwayO7.sigma0_cycleType}
\leanfile{ConwayO7/Canonical.lean}\leanok
\uses{def:sigma0}
\end{lemma}

\begin{proof}\leanok
\uses{lem:grp-cycletype-of-fixed, lem:grp-sigma0-order, lem:grp-sigma0-fixed-card}
Apply Proposition~\ref{lem:grp-cycletype-of-fixed} with $p = 7$, $f = 1$,
$m = 14$ and $99 = 1 + 7 \cdot 14$, using $\sigma_0^7 = 1$, $\sigma_0 \ne 1$
and the fixed-point count.
\end{proof}

\section{Transport and the reduction}

\begin{lemma}[Transport of strong regularity]\label{lem:grp-srg-transport}
If $G \simeq_g H$ is an isomorphism of finite graphs and $G$ is an
$\srg(n,k,\lambda,\mu)$, then so is $H$.
\lean{ConwayO7.isSRGWith_of_iso}
\leanfile{ConwayO7/Canonical.lean}\leanok
\uses{def:srg}
\end{lemma}

\begin{proof}\leanok
The isomorphism identifies neighbourhoods and common neighbourhoods
vertex-by-vertex, so all four parameters are preserved.
\end{proof}

\begin{theorem}[WLOG $\sigma_0$]\label{thm:grp-wlog}
If some $\srg(99,14,1,2)$ $\Gamma$, on any vertex set, admits an automorphism
of order $7$, then some $\srg(99,14,1,2)$ on the vertex set $\{0,\dots,98\}$
is invariant under $\sigma_0$.
\lean{ConwayO7.exists_sigma0_invariant}
\leanfile{ConwayO7/Canonical.lean}\leanok
\uses{def:srg, def:sigma0}
\end{theorem}

\begin{proof}[Proof sketch]\leanok
\uses{lem:grp-srg-transport, cor:grp-aut-cycle-type, lem:grp-sigma0-cycle-type}
Transport $\Gamma$ to $\{0,\dots,98\}$ along any bijection of vertex sets
(possible since $|V| = 99$); by Lemma~\ref{lem:grp-srg-transport} the
transported graph is again an $\srg(99,14,1,2)$, and the automorphism pushes
through to an order-seven adjacency-preserving permutation $\pi_1$.  By
Corollary~\ref{cor:grp-aut-cycle-type} and
Lemma~\ref{lem:grp-sigma0-cycle-type}, $\pi_1$ and $\sigma_0$ have the same
cycle type $[7^{14}]$, hence are conjugate in the symmetric group
(mathlib's \texttt{Equiv.Perm.isConj\_iff\_cycleType\_eq}); say
$\tau \pi_1 \tau^{-1} = \sigma_0$.  Relabelling once more by $\tau$ yields an
$\srg(99,14,1,2)$ invariant under $\sigma_0$ itself.
\end{proof}

\begin{theorem}[Reduction to $\sigma_0$]\label{thm:reduction}
Suppose no $\srg(99,14,1,2)$ on the vertex set $\{0,\dots,98\}$ is invariant
under $\sigma_0$.  Then $7 \nmid |\Aut\Gamma|$ for every $\srg(99,14,1,2)$
$\Gamma$.
\lean{ConwayO7.seven_not_dvd_card_aut_of_no_sigma0_invariant}
\leanfile{ConwayO7/Canonical.lean}\leanok
\uses{def:srg, def:sigma0}
\end{theorem}

\begin{proof}\leanok
\uses{lem:cauchy, thm:grp-wlog}
If $7 \mid |\Aut\Gamma|$ then by the Cauchy step (Lemma~\ref{lem:cauchy})
some automorphism has order $7$, and Theorem~\ref{thm:grp-wlog} produces a
$\sigma_0$-invariant $\srg(99,14,1,2)$ on $\{0,\dots,98\}$ --- contradicting
the hypothesis.  The hypothesis is exactly the statement encoded as a CNF
formula in Chapter~\ref{sec:encoding} and refuted by the verified proof
checker, completing this branch of Theorem~\ref{thm:main}.
\end{proof}

%=====================================================================
